\section*{Chapter \ref{ch:b1:clastic-sediments} --- Clastic Sediments and Their Environments}

\begin{solution}{exo:b1:clastic-sediments:1}
$\varphi = -\log_2 d$: \qty{0.18}{mm} gives $2.47$, a fine sand (2 to 3);
\qty{0.75}{mm} gives $0.42$, a coarse sand (0 to 1); \qty{3}{mm} gives
$-1.58$, a granule ($-1$ to $-2$); \qty{0.012}{mm} gives $6.38$, a fine silt (6
to 7); \qty{150}{mm} gives $-7.23$, a cobble ($-6$ to $-8$).
\end{solution}

\begin{solution}{exo:b1:clastic-sediments:2}
$M_z = (1.05 + 1.71 + 2.43)/3 = 1.73\,\varphi$, that is $2^{-1.73} =
\qty{0.30}{mm}$, a medium sand. $\sigma_I = (2.43 - 1.05)/4 + (2.98 - 0.62)/6.6
= 0.345 + 0.358 = 0.70\,\varphi$: moderately well sorted (0.50 to 0.71), just
below the limit of moderately sorted.
\end{solution}

\begin{solution}{exo:b1:clastic-sediments:3}
$q_b = (1 - 0.4) \times 0.02 \times (\num{0.01}/60)/2 = \qty{1.0e-6}{m^2.s^{-1}}$
of solid grains. In a day, $\num{1.0e-6} \times 86400 = \qty{0.0864}{m^3}$ per
metre of crest, or $0.0864 \times 2650 = \qty{229}{kg}$ of quartz.
\end{solution}

\begin{solution}{exo:b1:clastic-sediments:4}
(a)~50 knots are $50 \times 1852/3600 = \qty{25.7}{m.s^{-1}}$; 12 knots,
\qty{6.17}{m.s^{-1}}. (b)~$\qty{100000}{m}/\qty{6.17}{m.s^{-1}} =
\qty{16200}{s}$, about \qty{4.5}{h}. (c)~480 miles are \qty{772}{km}, 400 miles
\qty{644}{km}.
\end{solution}

\begin{solution}{exo:b1:clastic-sediments:5}
(a)~$(\varphi_{84} - \varphi_{16})/2 = 2 \times 0.9945\,\sigma/2 =
0.9945\,\sigma$, and $(\varphi_{95} - \varphi_5)/3.3 = 2 \times
1.6449\,\sigma/3.3 = 0.9969\,\sigma$: both estimate $\sigma$ to within a third
of a per cent. (b)~A real sediment need not be normal in phi units: a small
admixture of much coarser or much finer grains changes $\varphi_5$ or
$\varphi_{95}$ a lot and $\varphi_{16}$ and $\varphi_{84}$ little, so the second
estimate grows while the first does not. The two differ most for mixtures, such
as a sand with some gravel and mud, or a bimodal deposit.
\end{solution}

\begin{solution}{exo:b1:clastic-sediments:6}
(a)~$\sin\theta = T/L = 1.5/15 = 0.1$, so $\theta = \ang{5.74}$. (b)~$a =
c\tan\theta = 0.5 \times 0.1005 = \qty{0.050}{cm}$ a minute, or
\qty{3.0}{cm} an hour. (c)~A bed raised by centimetres an hour receives sand
falling out of suspension very fast: a heavily loaded current that is slowing
down, such as a waning flood or the tail of a \omterm{def:b1:clastic-sediments:turbidite}{turbidity current}.
\end{solution}

\begin{solution}{exo:b1:clastic-sediments:7}
(a)~$(350 + 20 + 5 + 340 + 15 + 30)/6 = 760/6 = 126.7$: a south-easterly
direction, although every reading lies within \ang{20} of north. (b)~$S =
\sum\sin\alpha_i = 0.672$, $C = \sum\cos\alpha_i = 5.692$, so $\bar\alpha =
\operatorname{atan2}(0.672, 5.692) = \ang{6.7}$, and $\bar R =
\sqrt{0.672^2 + 5.692^2}/6 = 5.732/6 = 0.955$. (c)~The arithmetic mean treats
350 and 005 as far apart, while on the circle they are \ang{15} apart: \omterm{def:b1:spherical-astronomy:horizontal}{azimuths}
must be averaged as unit vectors. The large $\bar R$ shows that the readings are
tightly grouped around \ang{7}.
\end{solution}

\begin{solution}{exo:b1:clastic-sediments:8}
A has more than \qty{5}{\percent} of clay: immature, and compositionally
immature too, with its \omterm{def:b1:rock-forming-minerals:groups}{feldspar} and rock fragments; quickly deposited near its
source, on an alluvial fan or a floodplain, or by a \omterm{def:b1:clastic-sediments:turbidite}{turbidity current}. B has less
than \qty{5}{\percent} clay and is well sorted ($\sigma_I < 0.5\,\varphi$) but its
grains are subangular ($\rho < 3$): mature, for example a river bar or a shelf
sand winnowed by currents. C has no clay, is very well sorted and its grains are
rounded ($\rho > 3$): supermature, and with \qty{98}{\percent} quartz also
compositionally mature, as on a beach or in a desert \omterm{def:b1:clastic-sediments:bedform}{dune} field.
\end{solution}

\begin{solution}{exo:b1:clastic-sediments:9}
(a)~$1.5/\sqrt{9.81 \times 0.3} = 0.87$; $1.0/\sqrt{9.81 \times 2} = 0.23$.
(b)~$h = U^2/g = 1.5^2/9.81 = \qty{0.23}{m}$; $U = \sqrt{9.81 \times 2} =
\qty{4.43}{m.s^{-1}}$. (c)~The Froude number grows as the depth decreases at a
given speed: thin, fast sheets of water, like the backwash of a wave on a beach,
easily reach Froude numbers near 1, while a river several metres deep would
need several metres per second, which few rivers reach except in floods.
\end{solution}

\begin{solution}{exo:b1:clastic-sediments:10}
(a)~With the mid-points $-0.5, 0.5, \ldots, 4.5$: $m = (-1 + 4 + 45 + 100 + 52.5
+ 22.5)/100 = 2.23\,\varphi$; the squared deviations weighted by the
percentages sum to \num{107.71}, so $s = \sqrt{1.0771} = 1.04\,\varphi$:
poorly sorted. (b)~The cumulative per cent coarser is 0, 2, 10, 40, 80, 95 and
100 at $-1, 0, \ldots, 5\,\varphi$. Then $\varphi_5 = 0 + 3/8 = 0.375$,
$\varphi_{16} = 1 + 6/30 = 1.20$, $\varphi_{50} = 2 + 10/40 = 2.25$,
$\varphi_{84} = 3 + 4/15 = 3.27$, $\varphi_{95} = 4.0$. $M_z = (1.20 + 2.25 +
3.27)/3 = 2.24\,\varphi$ and $\sigma_I = (3.27 - 1.20)/4 + (4.0 - 0.375)/6.6 =
0.517 + 0.549 = 1.07\,\varphi$. (c)~The two methods agree to about
\num{0.01} on the mean and \num{0.03} on the \omterm{def:b1:clastic-sediments:sorting}{sorting}: a fine sand
($2.24\,\varphi$ is \qty{0.21}{mm}), poorly sorted.
\end{solution}

\begin{solution}{exo:b1:clastic-sediments:11}
(a)~$U = \sqrt{0.02 \times 9.81 \times 50 \times 0.005/0.003} =
\qty{4.04}{m.s^{-1}}$, or 7.9 knots. (b)~$(\Delta\rho/\rho)\,h = U^2
C_f/(gS) = 25.7^2 \times 0.003/(9.81 \times 0.005) = \qty{40}{m}$: for instance
a current \qty{400}{m} thick with a density excess of \qty{10}{\percent}, which
shows how large and dense the current of 1929 must have been. (c)~The speed
scales as $\sqrt S$: $4.04/\sqrt 5 = \qty{1.81}{m.s^{-1}}$. The slower current
has weaker turbulence and drops its coarser grains; it becomes less dense, slows
further, and spreads its load over the plain as a \omterm{def:b1:clastic-sediments:turbidite}{graded bed}.
\end{solution}

\begin{solution}{exo:b1:clastic-sediments:12}
(a)~$0.7 \times \qty{5e10}{kg}/\qty{2000}{kg.m^{-3}} = \qty{1.75e7}{m^3}$ a
year. (b)~Spread over a front \qty{30}{m} high and \qty{100}{km} wide, it
advances the front by $\num{1.75e7}/(30 \times \num{1e5}) = \qty{5.8}{m}$ a
year. (c)~$a = v\tan\alpha = 5.83 \times \tan\ang{0.5} = \qty{0.051}{m}$, about
\qty{5}{cm} a year, a hundred times less than the advance; the
\qty{30}{m} succession forms in $30/0.051 = 590$ years.
\end{solution}

\begin{solution}{pb:b1:clastic-sediments:1}
\textbf{1.} Unit~1: $(5.0 + 6.6 + 8.4)/3 = 6.67\,\varphi$, \qty{0.0098}{mm}
(\qty{9.8}{\micro m}); unit~2: $3.47\,\varphi$, \qty{0.090}{mm}; unit~3:
$1.80\,\varphi$, \qty{0.29}{mm}; unit~4: $0.73\,\varphi$, \qty{0.60}{mm}.

\textbf{2.} Unit~1: $(8.4 - 5.0)/4 + (9.6 - 4.2)/6.6 = 0.850 + 0.818 =
1.67\,\varphi$; unit~2: $0.550 + 0.545 = 1.10$; unit~3: $0.275 + 0.273 = 0.55$;
unit~4: $0.725 + 0.712 = 1.44$.

\textbf{3.} Unit~1: fine silt, poorly sorted; unit~2: very fine sand, poorly
sorted; unit~3: medium sand, moderately well sorted; unit~4: coarse sand,
poorly sorted.

\textbf{4.} Unit~3. At the river mouth the current and the waves move the sand
again and again and carry the fines away, while the prodelta mud settles from
suspension as a mixture of all the fine sizes that reach it.

\textbf{5.} The size $8\,\varphi$ lies $1.4/1.8$ of the way from
$\varphi_{50}$ to $\varphi_{84}$, so $50 + (1.4/1.8) \times 34 = 76$ per cent of
the sample is coarser and about \qty{24}{\percent} is clay-sized. Such fine
particles settle only from nearly still water: the site was below the reach of
waves and currents.

\textbf{6.} A channel carries both gravel, rolled on its bed in floods, and
sand; the deposit mixes the gravel lag of the deepest flow with the sand of
lower flows, so it covers many classes.

\textbf{7.} The mean grows from \qty{0.0098}{mm} to \qty{0.090}{mm} and
\qty{0.29}{mm}: the succession coarsens upward.

\textbf{8.} Unit~1, laminated mud: the prodelta; unit~2, alternating silt and
sand: the lower delta front; unit~3, cross-bedded sand: the mouth bars of the
upper delta front.

\textbf{9.} Each zone came to lie on the one that was formerly in front of it:
the delta built forward, seaward, over the site.

\textbf{10.} At least \qty{21}{m}, the thickness of units 1 to 3, since the
succession filled the water from the floor of the prodelta to the mouth bar
near the surface; more, since the muds have been compacted since
(\cref{ch:b1:diagenesis-basins}).

\textbf{11.} A distributary channel of the delta plain, cutting into the mouth
bar as the delta built forward; its pebbles are its \omterm{def:b1:erosion-transport:load}{bed load}.

\textbf{12.} No: a point bar starts with an erosion surface and a lag and fines
upward, to floodplain mud. The delta front coarsens upward from mud.

\textbf{13.} $\qty{16}{m}/\num{4000}\text{ years} = \qty{4.0}{mm}$ a year.

\textbf{14.} $v = a/\tan\alpha = 0.0040/\tan\ang{0.5} = 0.0040/0.008727 =
\qty{0.458}{m}$ a year.

\textbf{15.} $0.458 \times 4000 = \qty{1830}{m}$, about \qty{1.8}{km}.

\textbf{16.} $\qty{21}{m}/\qty{4.0}{mm.yr^{-1}} = 5250$ years.

\textbf{17.} The delta front receives most of the sand near the mouths, the
prodelta far less and the delta plain little except in floods; the active mouths
also shift from place to place, so a point may receive much for a while and then
almost nothing. A single rate between two dated beds is an average over these
changes, and the muds have lost part of their thickness by compaction.

\textbf{18.} The sum is 1544, so the arithmetic mean is \ang{128.7}, a
south-easterly direction: meaningless, since the readings scatter around north
and the values near 340 and near 010 are averaged as if they were far apart.

\textbf{19.} $S = \sum\sin\alpha_i = 1.698$ and $C = \sum\cos\alpha_i =
11.248$.

\textbf{20.} $\bar R = \sqrt{1.698^2 + 11.248^2}/12 = 11.376/12 = 0.948$: the
readings are tightly grouped around one direction.

\textbf{21.} The rivers flowed, and the delta built, toward about north
(\cref{prop:b1:clastic-sediments:azimuths}); the land lay to the south.

\textbf{22.} Because the foresets were measured on beds that are no longer
horizontal: their \omterm{def:b1:geological-maps:attitude}{dip directions} are those of the tilted beds plus the
tilting. Rotating each foreset about the \omterm{def:b1:geological-maps:attitude}{strike} of the bedding by the \omterm{def:b1:geological-maps:attitude}{dip} of
the bedding restores the bedding, and with it the foresets, to their original
\omterm{def:b1:geological-maps:attitude}{attitude}.

\textbf{23.} Tidal currents flow in and out, and build sets that \omterm{def:b1:geological-maps:attitude}{dip} in two
opposite directions, as in herringbone \omterm{def:b1:clastic-sediments:cross-bedding}{cross-bedding}; the rose diagram would
have two lobes about \ang{180} apart. The twelve readings form one lobe,
$\bar R$ near 1: a river current flowing one way.

\textbf{24.} $\bar\alpha = \operatorname{atan2}(1.698, 11.248) = \ang{8.6}$: the
mean \omterm{def:b1:clastic-sediments:cross-bedding}{palaeocurrent} \omterm{def:b1:spherical-astronomy:horizontal}{azimuth} of unit~3 is \ang{9}, a little east of north.
\end{solution}
